\section{Related Work}
\label{sec:related}

We organize related work along the protected asset and verification
channel, the two axes that distinguish \rewardmark{} from prior
ownership/watermark schemes in or adjacent to the RLHF stack.

\paragraph{Watermarking the policy LLM's text outputs.}
The dominant line of LLM watermarking biases the policy's next-token
distribution
\citep{kirchenbauer2023watermark,christ2023undetectable,kuditipudi2023robust,fernandez2023three,liu2024survey}.
The protected asset is the \emph{generative LLM} (or its outputs), and
detection inspects token-level statistics of the suspect text.
\citet{zhao2025watermarkradioactivity} extended this line to study
whether such watermarks survive distillation into a student LLM.
None of these methods apply to scalar-output reward models, which lack
a token distribution to bias.

\paragraph{Watermarking the preference dataset.}
\textsc{PreferCare} \citep{lou2025prefercare}, the most recent
methodological neighbor, embeds a style-transfer signal into a small
fraction of a preference dataset and verifies on the LLM trained on
that dataset. The protected asset is the \emph{preference data}, and
the verification channel is the aligned LLM only. \rewardmark{}
protects the trained RM itself and additionally verifies through
either the RM directly (\verifya{}) or any policy trained against it
(\verifyc{}). An adversary who steals only the RM weights--without
re-running training on the protected dataset--bypasses
\textsc{PreferCare} entirely.

\paragraph{Watermarking the prompt or soft-prompt.}
\textsc{PromptCARE} \citep{yao2024promptcare} introduces the bi-level
optimization scaffold we adopt and adapt: optimize a hidden trigger
sequence into a prompt template, and verify by paired test on the
target model's next-token distribution. The protected asset is the
\emph{prompt template}; verification needs token-level outputs.
\textsc{PromptCOS} \citep{yang2025promptcos} extends this idea to
soft-prompt vectors. Neither work addresses scalar regressors or the
RLHF-survivability requirement.

\paragraph{Watermarking the encoder.}
\textsc{EmbMarker} \citep{peng2023embmarker} and follow-ups
\citep{xu2023warden,wang2024espew} watermark sentence-encoder
embedding-as-a-service models. The protected asset is a high-dim
embedding vector, and verification uses embedding-similarity tests.
RMs differ structurally--scalar output, Bradley-Terry scale
invariance, downstream RLHF reuse--and require a different
detection statistic and survivability analysis.

\paragraph{RM-level attacks (mechanism mirror).}
\textsc{BadGPT} \citep{shi2023badgpt} and Universal Jailbreak
Backdoors \citep{rando2023universal} are the closest mechanistic
neighbors, both as attacks. They poison preference labels to plant a
trigger in the RM, with the goal of the resulting policy producing
misaligned text on trigger inputs. \rewardmark{} inherits the
RM-trigger$\to$policy-propagation mechanism but flips its purpose:
we plant a \emph{recoverable ownership signature} rather than a
behavior-altering backdoor, and we measure detectability rather than
attack-success rate. Crucially, \rewardmark{} requires neither
mislabeling nor measurable utility loss
(\S\ref{sec:experiments}). \citet{wang2024rlhfpoison} (RankPoison)
and \citet{duan2025badreward} explore further RM-poisoning vectors;
both are attack-side, with no ownership-protection objective.

\paragraph{LLM fingerprinting and verifiable training.}
\textsc{HuRef} \citep{zeng2024huref} fingerprints LLM parameter
directions; \textsc{vTune} \citep{zhang2024vtune} watermarks via
injected fine-tuning data. Both protect the LLM or its training
data, not the RM. \rewardmark{} is complementary.

\paragraph{DPO theory: pair-margin vs.\ sampling-marginal.}
\citet{razin2025likelihood} formalize \emph{likelihood displacement}:
DPO can increase the implicit-reward margin while decreasing the
chosen-response log-probability.
\citet{feng20243d,liu2025misspec} characterize the resulting
distribution shift; \citet{pal2024smaug} propose DPOP as a
corrective. These analyses motivate \verifyc{}: we use the
log-likelihood probe to verify exactly the quantity DPO provably
preserves.

\paragraph{Position summary.}
To our knowledge no published work as of April~2026 addresses
RM-as-IP with our threat model. \rewardmark{} fills this gap.
